% =============================================================
% ① 【最优先加载】基础工具与核心设置
% =============================================================
% 说明：xcolor 必须尽早加载以防冲突；microtype 改善排版质量
\usepackage[table, svgnames]{xcolor}  % 颜色支持 (table支持表格彩色, svgnames支持更多颜色)
\usepackage{microtype}                % 微排版 (改进字距/断行)
\usepackage{setspace}                 % 行距设置宏包
\setlength{\baselineskip}{20pt}       % 设定正文行距为固定值 20磅

% =============================================================
% ② 字体设置与特殊符号
% =============================================================
\usepackage{pifont}                   % 特殊符号 (如带圈数字)

% 字体设置 (配合 main.tex 的 ctexart)
\setmainfont{Times New Roman}         % 英文/数字主字体
\setCJKmainfont[AutoFakeBold=true]{SimSun} % 中文主字体强制设为宋体

% =============================================================
% ③ 页面布局与页眉页脚
% =============================================================
% 3.1 页面边距设置
\usepackage[a4paper,
  top=2.5cm, bottom=2.5cm, left=2.5cm, right=2.5cm, % 页边距
  headsep=0.5cm,   % 页眉文字底部到正文顶部的距离
  headheight=15pt, % 页眉高度
  footskip=0.75cm, % 正文底部到页脚底部的距离
]{geometry}

% 3.2 页眉页脚设置
\usepackage{fancyhdr}
\fancypagestyle{reportstyle}{       % 定义 'reportstyle' 风格
  \fancyhf{} % 清空当前设置
  % 页眉：宋体小五，居中，显示章节标题
  \fancyhead[C]{\songti\zihao{-5} \leftmark}
  % 页脚：宋体小五，居中，显示页码
  \fancyfoot[C]{\songti\zihao{-5} \thepage}
  % 页眉横线
  \renewcommand{\headrulewidth}{0.5pt} 
}

% =============================================================
% ④ 数学公式与符号
% =============================================================
\usepackage{amsmath}    % 基础数学环境
\usepackage{amssymb}    % 更多数学符号
\usepackage{mathtools}  % amsmath 增强补丁 (需后加载)

% =============================================================
% ⑤ 图表与浮动体设置
% =============================================================
\usepackage{graphicx}   % 插图支持
\usepackage{float}      % 强行不浮动 [H]
\usepackage{subcaption} % 子图支持

% 5.1 表格支持
\usepackage{array}      % 基础表格扩展
\usepackage{tabularx}   % 自动列宽
\newcolumntype{C}{>{\centering\arraybackslash}X} % 定义居中自动列
\usepackage{booktabs}   % 三线表 (\toprule, \midrule, \bottomrule)
\usepackage{threeparttable} % 表格附注

% 5.2 题注格式 (Caption)
\usepackage{caption}
\DeclareCaptionFont{wuhao}{\zihao{5}\songti} % 定义宋体五号
\captionsetup{
    font=wuhao,
    labelsep=quad,         % 序号后空一汉字宽
    justification=centering,
    singlelinecheck=true
}
% 表格标题特殊设置 (标题在上方)
\captionsetup[table]{position=top, skip=6pt}
% 图标题特殊设置 (标题在下方)
\captionsetup[figure]{position=bottom}

% 5.3 图表编号格式 (图 2-1)
\numberwithin{figure}{section}
\numberwithin{table}{section}
\renewcommand{\thefigure}{\thesection-\arabic{figure}}
\renewcommand{\thetable}{\thesection-\arabic{table}}

% =============================================================
% ⑥ 功能扩展 (列表、代码、URL)
% =============================================================
\usepackage{enumitem}   % 列表定制
\usepackage{xurl}       % URL 自动换行 (优化参考文献排版)

% 代码高亮设置
\usepackage{listings}
\renewcommand{\lstlistingname}{代码}
\lstdefinestyle{cstyle}{
  language=C,
  basicstyle=\ttfamily\small,
  numbers=left,
  numberstyle=\scriptsize\color{gray},
  numbersep=8pt,
  commentstyle=\color{gray!70!black},
  showstringspaces=false,
  tabsize=4,
  frame=single,
  rulecolor=\color{teal!70!blue!30},
  frameround=tttt,
  columns=fullflexible,
  keepspaces=true
}
\lstset{style=cstyle}

% =============================================================
% ⑦ 章节标题与目录样式 (Styling)
% =============================================================
% 7.1 章节标题格式 (ctexset)
\ctexset{
    section = {
        format = \zihao{3}\heiti\centering,
        beforeskip = 30pt,
        afterskip = 20pt,
        number = \arabic{section}
    },
    subsection = {
        format = \zihao{4}\heiti\raggedright,
        beforeskip = 25pt,
        afterskip = 12pt,
    },
    subsubsection = {
        format = \zihao{-4}\heiti\raggedright,
        beforeskip = 12pt,
        afterskip = 6pt,
    },
    contentsname = {目\hspace{2em}录} % 目录标题格式
}
\setcounter{secnumdepth}{3} % 编号深度

% =============================================================

% 7.2 目录样式 (titletoc)
\usepackage{titletoc}
\setcounter{tocdepth}{2}    % 目录显示深度

% 一级标题 (Section)
\titlecontents{section}[0pt]
    {\addvspace{0.4em}\heiti\zihao{-4}}
    {\thecontentslabel\hspace{1em}}
    {}
    {\titlerule*[0.6pc]{.}\contentspage}

% 二级标题 (Subsection)
\titlecontents{subsection}[2em]
    {\addvspace{0.1em}\songti\zihao{-4}}
    {\thecontentslabel\hspace{1em}}
    {}
    {\titlerule*[0.6pc]{.}\contentspage}

% 三级标题 (Subsubsection)
\titlecontents{subsubsection}[4em]
    {\addvspace{0em}\songti\zihao{-4}}
    {\thecontentslabel\hspace{1em}}
    {}
    {\titlerule*[0.6pc]{.}\contentspage}

% =============================================================
% ⑧ 参考文献设置
% =============================================================
\usepackage[style=gb7714-2015, backend=biber, sorting=none]{biblatex}
\addbibresource{bib/ref.bib}

% 外观定制
\renewcommand{\bibfont}{\zihao{5}\songti} % 宋体五号
\setlength{\bibitemsep}{3pt}              % 条目间距
\setlength{\bibnamesep}{0pt}
\setlength{\bibinitsep}{0pt}
\setlength{\biblabelsep}{1.5em}           % 序号与内容间距 (悬挂缩进)

% URL 断行优化
\setcounter{biburllcpenalty}{7000}
\setcounter{biburlucpenalty}{8000}

% 标题间距调整
\defbibheading{bibliography}[\bibname]{%
  \section*{#1}%
  \markboth{#1}{#1}%
  \addvspace{20pt}
}

% =============================================================
% ⑨ 【最后加载】超链接与智能引用
% =============================================================
\usepackage{hyperref}
\hypersetup{
  colorlinks=false,   % 不用彩色文字
  pdfborder={0 0 0}   % 去掉红框
}

\usepackage{cleveref} % 必须在 hyperref 之后
\crefname{lstlisting}{代码}{代码}
\providecommand{\lstlistingautorefname}{代码}
\renewcommand{\figurename}{图}
\renewcommand{\figureautorefname}{图}
\renewcommand{\tablename}{表}
\renewcommand{\tableautorefname}{表}
\renewcommand{\sectionautorefname}{章}
\renewcommand{\subsectionautorefname}{节}
\renewcommand{\subsubsectionautorefname}{小节}



